\section{Total Derivatives}

\subsection{Type, Domain, and Codomain}

\begin{enumerate}
  \item Let $f : \mathbb{R}^n \to \mathbb{R}^m$. What is the type (domain and codomain) of the total derivative $Df(\mathbf{x})$ at a point $\mathbf{x} \in \mathbb{R}^n$?
  \item How can you interpret $Df(\mathbf{x})$ as a linear map between tangent spaces?
  \item In coordinates, how is $Df(\mathbf{x})$ represented as a matrix? What are its dimensions?
\end{enumerate}

\subsection{Jacobian as a Limit (First Principles)}

\begin{enumerate}[resume]
  \item Starting from the definition of differentiability, write the limit definition of $Df(\mathbf{x})$.
  \item Show that if $f$ is differentiable at $\mathbf{x}$, then there exists a matrix $J_f(\mathbf{x})$ such that
    \[
      \lim_{\|\mathbf{h}\| \to 0} \frac{\|f(\mathbf{x}+\mathbf{h}) - f(\mathbf{x}) - J_f(\mathbf{x})\mathbf{h}\|}{\|\mathbf{h}\|} = 0.
    \]
  \item How do the entries of $J_f(\mathbf{x})$ arise from partial derivatives?
  \item Why does this limit definition guarantee that the Jacobian gives the \emph{best linear approximation}?
\end{enumerate}

\subsection{Directional Derivative and Gradient}

\begin{enumerate}[resume]
  \item How does the total derivative differ from a partial derivative?
  \item How does the total derivative relate to the gradient $\nabla f$?
  \item What happens to the total derivative if one of the variables is held constant?
  \item How is the total derivative related to directional derivatives?
  \item Define the directional derivative of $f$ at $\mathbf{x}$ in the direction $\mathbf{v}$ using a limit.
  \item Show that the directional derivative can be expressed as
    \[
      Df(\mathbf{x})[\mathbf{v}] = J_f(\mathbf{x}) \mathbf{v}.
    \]
  \item In the scalar-valued case $f : \mathbb{R}^n \to \mathbb{R}$, how is the Jacobian related to the gradient $\nabla f(\mathbf{x})$?
  \item Why does the gradient point in the direction of steepest ascent?
\end{enumerate}

\subsection{Tangent Planes and Local Linearization}

\begin{enumerate}[resume]
  \item For $f : \mathbb{R}^n \to \mathbb{R}$, write the equation of the tangent hyperplane at $\mathbf{x}_0$.
  \item How does the linear map $Df(\mathbf{x}_0)$ define this hyperplane?
  \item What is the first-order Taylor approximation of $f$ near $\mathbf{x}_0$?
  \item Why is $Df(\mathbf{x}_0)$ called the \emph{local linearization} of $f$?
\end{enumerate}

\subsection{Normal Vector}

\begin{enumerate}[resume]
  \item For a level set $f(\mathbf{x}) = c$, what vector is normal to the surface at $\mathbf{x}$?
  \item How is this normal vector expressed in terms of the Jacobian?
  \item In the case $f : \mathbb{R}^2 \to \mathbb{R}$, how does $\nabla f$ relate to the tangent line?
\end{enumerate}

\subsection{Linear Map and Best Linear Approximation}

\begin{enumerate}[resume]
  \item What does it mean for $Df(\mathbf{x})$ to be a linear map?
  \item What exactly is being mapped by $Df(\mathbf{x})$?
  \item Why is $Df(\mathbf{x})$ the \emph{best} linear approximation to $f$ near $\mathbf{x}$?
  \item How does the error behave compared to $\|\mathbf{h}\|$?
\end{enumerate}

\subsection{Jacobian Determinant and Change of Variables}

\begin{enumerate}[resume]
  \item For $f : \mathbb{R}^n \to \mathbb{R}^n$, what is the Jacobian determinant $\det J_f(\mathbf{x})$?
  \item How does $\det J_f(\mathbf{x})$ describe local volume (or area) scaling?
  \item Why does $|\det J_f(\mathbf{x})|$ appear in the change of variables formula for multiple integrals?
  \item What happens geometrically when $\det J_f(\mathbf{x}) = 0$?
\end{enumerate}

\subsection{Multivariable Chain Rule}

\begin{enumerate}[resume]
  \item Let $f : \mathbb{R}^n \to \mathbb{R}^m$ and $g : \mathbb{R}^m \to \mathbb{R}^k$. State the chain rule for $D(g \circ f)(\mathbf{x})$.
  \item Show that
    \[
      J_{g \circ f}(\mathbf{x}) = J_g(f(\mathbf{x})) \, J_f(\mathbf{x}).
    \]
  \item How does this relate to composition of linear maps?
\end{enumerate}

\subsection{Connection to Taylor Series}

\begin{enumerate}[resume]
  \item Write the first-order Taylor expansion of $f : \mathbb{R}^n \to \mathbb{R}^m$ at $\mathbf{x}_0$.
  \item Where does the Jacobian appear in this expansion?
  \item How would higher-order derivatives extend this approximation?
\end{enumerate}

\subsection{Newton's Method for Systems}

\begin{enumerate}[resume]
  \item Let $F : \mathbb{R}^n \to \mathbb{R}^n$. Write the Newton iteration for solving $F(\mathbf{x}) = \mathbf{0}$.
  \item How is the Jacobian matrix used in each iteration?
  \item Why does Newton's method rely on local linearization?
  \item What conditions on $J_F(\mathbf{x})$ are required for the method to work?
\end{enumerate}

\subsection{Mapping Tangent Vectors}

\begin{enumerate}[resume]
  \item Let $f : \mathbb{R}^n \to \mathbb{R}^m$ and let $\gamma : \mathbb{R} \to \mathbb{R}^n$ be a differentiable curve with $\gamma(0) = x_0$. What is the type (domain and codomain) of $\gamma'(0)$ and of $Df(x_0),\gamma'(0)$?
  \item Given a curve $\gamma(t)$ with $\gamma(0) = x_0$, write the expression for the tangent vector to the image curve $f(\gamma(t))$ at $t=0$. Which theorem justifies this expression?
  \item State and explain the identity
    \[
      \frac{d}{dt} f(\gamma(t)) \bigg|_{t=0} = Df(x_0)\,\gamma'(0).
    \]
    What are the roles of each term in this equation?
  \item Why does the Jacobian $Df(x_0)$ depend only on the tangent vector $\gamma'(0)$ and not on the specific curve $\gamma$? What does this imply about its geometric meaning?
  \item Suppose two curves $\gamma_1(t)$ and $\gamma_2(t)$ satisfy $\gamma_1(0)=\gamma_2(0)=x_0$ and $\gamma_1'(0)=\gamma_2'(0)$. What can you say about the derivatives of their image curves under $f$ at $t=0$?
  \item Show that for any vector $v \in \mathbb{R}^n$, the curve $\gamma(t) = x_0 + t v$ satisfies
    \[
      Df(x_0)v = \frac{d}{dt} f(x_0 + t v)\bigg|_{t=0}.
    \]
    What does this tell you about the relationship between the Jacobian and directional derivatives?
  \item In what sense can the Jacobian $Df(x_0)$ be viewed as a map between tangent spaces? Write explicitly the mapping it defines.
  \item Explain what is meant by the statement: ``the Jacobian maps tangent vectors at $x_0$ to tangent vectors at $f(x_0)$.'' What are the spaces involved?
  \item Consider all curves passing through $x_0$. What structure do their tangent vectors at $t=0$ form? How does the Jacobian act on this structure?
  \item Why is the Jacobian called the ``best linear approximation'' of $f$ at $x_0$? How does the curve-based interpretation support this claim?
  \item Let $f : \mathbb{R}^2 \to \mathbb{R}^2$ and let $\gamma(t)$ be a curve through $x_0$. Describe geometrically how the tangent vector $\gamma'(0)$ is transformed by $Df(x_0)$.
  \item How does the mapping of tangent vectors by $Df(x_0)$ lead to the notion of a tangent line, plane, or hyperplane to the image of $f$?
  \item If $f : \mathbb{R}^n \to \mathbb{R}$, how does $Df(x_0)$ relate to the gradient $\nabla f(x_0)$? How does this affect the interpretation of $Df(x_0)v$?
  \item Explain why the Jacobian provides a linear map that approximates how small displacements near $x_0$ are transformed under $f$.
\end{enumerate}
